%
% Section: Collision and Animation in Crowd Simulation
% \section{Collision and Animation in Crowd Simulation}
\subsection{Scalable Interaction and Solver Frameworks}
\label{sec:solvers}

This section reviews key methodologies and recent advances in collision handling and animation strategies for large-scale crowd simulation, with emphasis on approaches that enable real-time performance for dense populations. The discussion synthesizes constraint-based positional methods, velocity-space collision avoidance, perceptual (foveated) animation optimizations, and hierarchical spatial partitioning frameworks that together address the computational bottlenecks of collision dynamics~\cite{weiss,charlton,stancu,sung}.

\subsection{Position-Based Dynamics (PBD)}
Position-Based Dynamics (PBD) has emerged as a robust alternative to explicit force-based schemes for dense crowds, because it enforces non-penetration and other constraints directly in the positional domain rather than through accumulated forces. By iteratively projecting positions to satisfy constraints, PBD mitigates numerical instabilities (e.g., jitter) that commonly arise with large timesteps or high agent densities. Implementations typically combine PBD constraint projection with a lightweight locomotion controller: agent velocities are blended with preferred velocities using an interpolation factor $\xi$ to produce smooth accelerations,
\[
v \gets v(1 - \xi) + \xi v_{\mathrm{pref}},
\]
and neighbor discovery is often implemented with spatial hashing to achieve near-$O(1)$ neighbor queries in practice~\cite{weiss}. PBD's direct position correction provides predictable convergence in high-contact regimes and scales well to tens of thousands of agents when paired with GPU-friendly neighborhood search structures~\cite{weiss}.

\subsection{Velocity-Space Collision Avoidance}
Velocity-space methods, most notably Optimal Reciprocal Collision Avoidance (ORCA), formulate per-agent constraints in velocity space to guarantee (under certain assumptions) locally collision-free motion. The core difficulty for large populations lies in solver scalability: ORCA requires solving small linear programs per agent, and naïve CPU implementations become the $O(n^{2})$ bottleneck as population counts grow. Recent work has focused on migrating ORCA solvers to GPUs using batch linear-program (LP) strategies and careful thread-workload balancing to preserve warp coherency. Grid-based spatial partitioning replaces tree structures in such GPU implementations to provide SIMD-friendly parallelism, memory coherence, and localized queries, producing large speedups (reported up to $\approx 30\times$ compared to optimized multi-core CPU baselines) and enabling interactive update rates for hundreds of thousands of agents under constrained visual fidelity trade-offs~\cite{charlton}.

\subsection{Perceptual Optimization: Foveated Animations}
Animation fidelity often dominates runtime cost when rendering and updating large crowds. Foveated animation leverages human visual limitations—specifically reduced peripheral acuity and the crowding effect—to allocate high-fidelity skeletal updates to the foveal region and progressively coarser representations in the periphery. Typical schemes partition visual eccentricity into zones (e.g., central fovea: full skeletal updates; immediate periphery: reduced update rates but above Critical Flicker Frequency (CFF); deep periphery: statistical or impostor motion only) and adapt update rates to maintain perceived continuity while drastically reducing per-agent animation cost. Empirical prototypes have reported very large reductions in animation operations (orders of magnitude in extreme cases) with minimal perceptual loss when thresholds are chosen based on CFF and eccentricity-dependent perceptual studies~\cite{stancu}.

\subsection{Hierarchical Spatial Partitioning and Rigid-Body Response}
For rigid-body interactions and narrow-phase contact resolution, hierarchical spatial partitioning (e.g., octree combined with AABB pruning) provides an effective defense against quadratic contact-check explosion. A shallow, GPU-traversal-friendly hierarchy—often a two-level scheme—strikes a balance between culling efficiency and traversal overhead. Collision response commonly employs a two-stage strategy: mass-weighted penetration correction (frequently using Baumgarte-style stabilization) to remove interpenetration reliably, followed by impulse-based velocity updates using Coulomb friction models to produce plausible restitution and tangential sliding or sticking behaviors. Such hierarchical frameworks have demonstrated robustness for complex meshes and have enabled interactive rigid-body simulations at scales where flat broad-phase methods fail~\cite{sung}.

%\subsection{Synthesis and Practical Considerations}
%Modern, large-scale crowd systems tend to combine multiple techniques: PBD or constraint projection for close-contact stability, GPU-accelerated velocity-space solvers for locomotion-level collision avoidance, foveated or perceptual LODs for animation, and hierarchical spatial structures for rigid-body culling. Practical system design requires careful profiling of (1) neighbor discovery and spatial data structure traversal, (2) per-agent solver workload balance on the GPU (to avoid warp divergence), and (3) perceptual thresholds that determine where fidelity can be reduced without visible artifacts. The combined approach allows interactive simulation for very large populations provided engineers accept calibrated perceptual and visual trade-offs~\cite{weiss,charlton,stancu,sung}.

\subsection{Geometric Representations of Agents}
\label{sec:geom}

\subsubsection{Disc-Based Models}
Discs are computationally cheap and isotropic; they simplify many algorithms but omit orientation and rotational behavior. The rotational symmetry that makes discs easy to compute is also their limiting factor in narrow passages and close-proximity interactions.

\subsubsection{Ellipse and EOBB Representations}
Ellipses and oriented bounding boxes introduce orientation and better ground-plane approximation at modest extra cost. They improve corridor behavior and reduce conservative collision footprints, but they complicate distance computations and projection formulas.

\subsubsection{Capsule-Based (Torso) Representations multi-agent torso crowds}
\label{sec:torso-crowds}
A recent thread of work shows that explicit torso modeling (2D capsules) captures per-agent rotation, shoulder clearance, and head profile in a simple geometry that aligns closely with human torso footprints~\cite{talukdar2025,stuevel2017}. Talukdar and Weiss (2025) present a PBD-based formulation with three interacting constraint families that directly address close-proximity collisions:

\begin{itemize}
  \item \textbf{Short-range interaction (contact constraint).} Compute the shortest point-pair between two capsule line segments and enforce a non-penetration constraint on the corresponding disc radii extracted from each capsule cross-section. If $\tilde{x}_i,\tilde{x}_j$ are the closest points, the inequality constraint is
  \[
  C(\tilde{x}_i,\tilde{x}_j) = \|\tilde{x}_i - \tilde{x}_j\| - (r_i + r_j) \ge 0,
  \]
  and violation is corrected by the standard PBD mass-weighted projection (see Eq.~(4) in \cite{talukdar2025}). This disc-superposition trick simplifies narrow-phase corrections while preserving the capsule geometry's benefits.
  \item \textbf{Personal-space constraint.} A soft spacing constraint increases comfort distance and invites early maneuvers. Talukdar et al. model it with a per-agent threshold (user-tunable) so that $C(\tilde{x}_i,\tilde{x}_j) \ge d_{\mathrm{personal}}$ and correct violations with reduced stiffness for organic, anticipatory motion.
  \item \textbf{Long-range (anticipatory) interaction.} Using predicted positions along each agent's velocity, future time-to-contact $g$ is computed (as in Weiss et al.), and if $g$ exceeds a configured threshold the algorithm applies softened, timed positional corrections scaled by a stiffness factor $\exp(-\hat{g}^2/g_0)$ to avoid sudden maneuvers. Talukdar et al. adapt disc-based future-collision calculations to capsule pairs using the closest-point projection described above, enabling anticipatory rotation and lane-forming behavior in tight spaces~\cite{talukdar2025}.
\end{itemize}

Beyond constraints, the torso work adds a \textbf{rotation planner} that computes available clearance along the agent's forward axis and derives a target rotation $\theta^\star$ (clamped per time-step) so agents can turn their torso to exploit local clearance (see Fig.~5 in \cite{talukdar2025}). The practical payoff is significant: capsules can pass where discs deadlock, and torso-aware PBD supports smoother close-proximity behavior while remaining amenable to PBD's constraint projection pipeline.
